%%
%% This is file `main.tex' based on `sample-sigconf.tex'
%%
\documentclass[sigconf]{acmart}

\usepackage[utf8]{inputenc}
\usepackage[portuguese]{babel}
\usepackage{listings}
\usepackage{color}
\usepackage{amsmath}

\definecolor{dkgreen}{rgb}{0,0.6,0}
\definecolor{gray}{rgb}{0.5,0.5,0.5}
\definecolor{mauve}{rgb}{0.58,0,0.82}

\lstset{frame=tb,
  language=Python,
  aboveskip=3mm,
  belowskip=3mm,
  showstringspaces=false,
  columns=flexible,
  basicstyle={\small\ttfamily},
  numbers=none,
  numberstyle=\tiny\color{gray},
  keywordstyle=\color{blue},
  commentstyle=\color{dkgreen},
  stringstyle=\color{mauve},
  breaklines=true,
  breakatwhitespace=true,
  tabsize=3
}

%%
%% \BibTeX command to typeset BibTeX logo in the docs
\AtBeginDocument{%
  \providecommand\BibTeX{{%
    \normalfont B\kern-0.5em{\scshape i\kern-0.25em b}\kern-0.8em\TeX}}}

%% Rights management information.
\setcopyright{acmlicensed}
\copyrightyear{2026}
\acmYear{2026}
\acmDOI{XXXXXXX.XXXXXXX}

%% These commands are for a PROCEEDINGS abstract or paper.
\acmConference[ESI 2026]{Engenharia de Software Inteligente - Mestrado UFU}{Julho 2026}{Uberlândia, MG}
\acmISBN{978-1-4503-XXXX-X/26/07}

\begin{document}

\title{Avaliação de LLMs como Assistentes de Codificação e Design Arquitetural: Um Estudo de Caso na Construção da Plataforma PLURI}

\author{Gabriel Camargos}
\email{gabriel.camargos@ufu.br}
\affiliation{%
  \institution{Universidade Federal de Uberlândia (UFU)}
  \city{Uberlândia}
  \state{Minas Gerais}
  \country{Brasil}
}

\begin{abstract}
O desenvolvimento de pipelines de processamento de documentos e integrações sistêmicas do zero exige um alto esforço cognitivo e técnico. Este trabalho investiga o papel dos Large Language Models (LLMs) como assistentes ativos no ciclo de vida de desenvolvimento de software (SDLC), focando no design arquitetural e na automação de pipelines. O estudo de caso descreve a construção do "Agente de Questões" para a plataforma de gestão educacional PLURI. O sistema automatiza a ingestão de um banco de questões legado (arquivos PDF e DOCX organizados por ano e bimestre) realizando a extração estruturada de dados via Pydantic. O agente identifica dinamicamente elementos multimídia nos documentos brutas, gera marcadores de imagem visualmente ordenados (\lstinline{[IMAGE_X]}), faz o upload automático para um servidor de armazenamento via API REST, e realiza a substituição por tags HTML \lstinline{<img>} nos campos respectivos do JSON de cadastro. Para avaliar a robustez do software co-autorado por IA, implementamos uma suíte completa de testes unitários e de integração (utilizando FastAPI e Pytest). Além disso, analisamos o trade-off de precisão, latência e custo financeiro (tokens consumidos) entre os modelos Gemini 1.5 e 2.0/2.5. Nossos resultados indicam que a integração de esquemas de dados estruturados com prompts dinâmicos reduz significativamente a taxa de alucinação e viabiliza a automação robusta de ponta a ponta.
\end{abstract}

\keywords{Large Language Models, Engenharia de Software Assistida, Design Arquitetural, Processamento Multimodal, Ingestão de Dados, API REST, Pytest, Gemini 2.5}

\maketitle

\section{Introdução}
A evolução dos Large Language Models (LLMs) tem provocado uma mudança profunda na Engenharia de Software. Estudos clássicos como o de Chen et al. \cite{chen2021evaluating} com a introdução do benchmark HumanEval demonstraram o poder dos modelos na geração de código funcional em nível de método. Mais recentemente, investigações como as de Peng et al. \cite{peng2023impact} documentaram ganhos empíricos significativos na produtividade dos desenvolvedores na escrita de códigos rotineiros utilizando assistentes inteligentes. 

Contudo, a Engenharia de Software moderna envolve desafios que transcendem a codificação de algoritmos isolados. Levantamentos abrangentes como os de Fan et al. \cite{fan2023large} e Hou et al. \cite{hou2024survey} apontam que os maiores desafios técnicos atuais residem no ciclo de vida completo do software (SDLC), englobando design arquitetural, modularidade, e integrações complexas. A transição da geração de código simples para a inferência e modelagem de arquiteturas baseadas em requisitos de linguagem natural é uma fronteira em desenvolvimento, conforme explorado por estudos de 2026 \cite{canai2026build, empirical2026architectural}. A IA enfrenta barreiras significativas ao tentar separar detalhes finos de implementação da abstração necessária para a visualização arquitetural de sistemas inteiros \cite{automated2026view}.

Este artigo aborda esse problema por meio de um estudo de caso prático: a concepção e implementação do \textbf{Agente de Ingestão de Questões} da plataforma educacional PLURI. O sistema precisava resolver um problema comum em sistemas corporativos: migrar um banco de dados legado de avaliações escolares (PDF e DOCX organizados em árvores de diretórios por ano e bimestre) para uma API REST estruturada. Mais do que extrair texto, o agente precisa realizar tarefas de design: identificar metadados contextuais, correlacionar disciplinas e assuntos com o banco existente e, crucialmente, isolar imagens das questões em sua ordem visual de leitura, fazer o upload para o storage remoto, obter as URLs públicas e formatar o corpo da questão em HTML com as tags \lstinline{<img>} correspondentes.

Para validar esta implementação, estruturamos uma suíte rigorosa de testes unitários e de integração utilizando Pytest e FastAPI, garantindo a conformidade dos contratos de dados modelados com Pydantic. Também detalhamos a análise de custos do processamento via LLMs para demonstrar a viabilidade econômica do sistema em produção.

\section{Fundamentação Teórica}

\subsection{LLMs no Desenvolvimento de Software}
Os LLMs têm demonstrado grande utilidade na automação de tarefas do SDLC \cite{fan2023large, hou2024survey}. Enquanto a geração de código out-of-the-box foca na conversão direta de especificações textuais em algoritmos funcionais \cite{chen2021evaluating}, a engenharia de software inteligente exige que estes modelos colaborem na estruturação de padrões arquiteturais e interfaces de API estáveis \cite{empirical2026architectural}.

\subsection{Design Arquitetural e Abstração com LLMs}
O planejamento estrutural de software com IA envolve a capacidade do modelo de compreender regras de domínio complexas. Estudos empíricos de 2026 \cite{canai2026build} sugerem que, embora LLMs consigam propor arquiteturas de alto nível de maneira satisfatória, sua precisão e consistência caem à medida que o escopo e o número de restrições do sistema aumentam. A geração automatizada de visões arquiteturais a partir de repositórios reais ainda sofre com a incapacidade da IA de isolar ruídos de implementação da semântica conceitual \cite{automated2026view}. Neste estudo, mitigamos esse problema utilizando o Pydantic como âncora contratual rígida que força o modelo a gerar saídas estruturadas estritas.

\subsection{O Papel de Schemas de Validação}
O Pydantic atua como uma interface de validação em tempo de execução que define os tipos e restrições dos dados de entrada. Na integração com LLMs, a especificação do schema Pydantic é injetada dinamicamente no prompt do modelo. Isso atua como um gabarito formal de saída, reduzindo drasticamente erros de desserialização JSON e garantindo compatibilidade contínua com endpoints REST.

\section{Arquitetura e Implementação do Agente}
A arquitetura do Agente PLURI foi estruturada para suportar o fluxo assíncrono de importação e processamento multimídia. O sistema é dividido em componentes modulares integrados:

\subsection{Parsing de Documentos com Ordem Visual e Desduplicação}
O módulo \lstinline{document_parser.py} foi atualizado para preservar a integridade estrutural dos documentos originais e otimizar a extração multimídia. A extração padrão de texto em PDFs frequentemente ignora a relação de proximidade espacial entre textos e imagens.
Utilizando a biblioteca PyMuPDF (\lstinline{fitz}), passamos a recuperar a estrutura do documento via \lstinline{page.get_text("dict")}. Isso permite ler blocos de conteúdo ordenados sequencialmente. Caso o bloco seja de texto, ele é concatenado. Caso seja de imagem, os bytes brutos são extraídos.

Para otimizar os recursos do storage do servidor, o \lstinline{document_parser} calcula a assinatura de hash MD5 de cada imagem em tempo real. Se o hash já tiver sido processado no mesmo documento (situação comum com logotipos escolares ou marca d'águas em provas legadas), o parser descarta a imagem duplicada e insere o marcador existente (e.g. \lstinline{[IMAGE_X]}) na posição textual correspondente. Caso contrário, um novo marcador de imagem é associado à fila de mídias. Uma lógica similar foi implementada para DOCX via análise XML das tags \lstinline{w:drawing} e \lstinline{a:blip}.

\subsection{Processamento e Upload de Mídia (image\_handler.py)}
O serviço de tratamento de imagens, \lstinline{image_handler.py}, processa as questões extraídas pela IA no seguinte fluxo:
1. Varre os campos \lstinline{corpo}, \lstinline{introducaoAlternativa} e o texto de cada alternativa procurando marcadores do tipo \lstinline{\[IMAGE_(\d+)\]}.
2. Para cada índice encontrado, consulta um cache local. Havendo uma falha de cache (miss), o agente faz uma requisição HTTP POST Multipart de upload de bytes para a rota \lstinline{/controle-de-arquivos/enviar/} do backend.
3. A API Spring Boot salva o arquivo e retorna uma URL pública.
4. O placeholder \lstinline{[IMAGE_X]} no texto é substituído por uma tag de imagem HTML:
   \lstinline{<p><img src="URL_PUBLICA" style="max-width: 100%; height: auto;"></p>}.
5. O objeto \lstinline{FileModel} contendo metadados é criado e associado à questão ou alternativa.

\subsection{Pipeline de Ingestão em Lote e Detecção de Metadados}
Criamos a rota assíncrona \lstinline{/ingest-question-bank} para automatizar a ingestão em larga escala. Ela varre recursivamente a pasta \lstinline{data/banco_questoes}, estruturada sob o padrão \lstinline{/data/banco_questoes/[ANO]-[PERIODO]/[ARQUIVO_PROVA].pdf}. O agente extrai o \lstinline{ano} (e.g., \lstinline{2023}) diretamente da pasta pai. Os metadados de disciplinas e assuntos são buscados via API REST com base na área detectada dinamicamente pelo LLM (Processamento em Duas Fases). O agente injeta o valor de \lstinline{ano} no objeto Pydantic de cada questão. O andamento do processo é gravado em tempo real em um arquivo markdown de auditoria (\lstinline{logs/ingestion_log.md}) e pode ser consultado a qualquer instante através da rota \lstinline{/ingestion-status}.

\subsection{Mitigação de Limites de Cota de API por Chunking com Janela Deslizante}
Ao processar arquivos PDF extensos (e.g., acima de 25.000 caracteres), a quantidade excessiva de tokens submetida à API do LLM em uma única requisição estourava com facilidade a quota de TPM (Tokens por Minuto) da Vertex AI, resultando em erros HTTP `429 Resource Exhausted`. 

Para mitigar este gargalo, implementamos um algoritmo de processamento por janelas deslizantes (Sliding-Window Chunking). O texto completo da prova é fatiado em blocos (chunks) de 20.000 caracteres com uma sobreposição (overlap) de 3.000 caracteres para evitar que alguma questão seja truncada no corte. Introduzimos um atraso (sleep) assíncrono de 2 segundos entre os blocos, permitindo que a janela de quota se recupere. Para evitar duplicidades causadas pela sobreposição de blocos, o agente executa uma limpeza espacial pós-extração, de-duplicando as questões por meio de chaves unificadas e normalizadas dos títulos e corpos de texto antes de salvá-las ou cadastrá-las.

\subsection{Justificativa e Fundamentação da Arquitetura de Pipeline}
A arquitetura do agente adota o padrão arquitetural Pipes and Filters (Canos e Filtros) \cite{garlan1993introduction}. Ao invés de submeter o documento bruto diretamente a um modelo de IA multimodal em uma única etapa (o que acarretaria em altíssimos custos de processamento, estouro de contexto e baixa confiabilidade na ordenação espacial de leitura), o pipeline separa as tarefas de forma determinística:
\begin{itemize}
\item \textbf{Filtro de Extração Textual e Layout:} O parsing de baixo nível (PyMuPDF) isola a complexidade estrutural, preservando posições relativas e convertendo fluxos de bytes de imagens em marcadores unificados.
\item \textbf{Filtro de Desduplicação por Assinatura:} O algoritmo calcula hashes MD5 para evitar processar mídias repetidas.
\item \textbf{Filtro de Orquestração Semântica (LLM):} O LLM atua sob demanda focando apenas no texto delimitado por um schema rigoroso Pydantic \cite{pydantic2024}, limitando as chances de alucinações sintáticas.
\item \textbf{Filtro de Processamento Multimídia (Upload):} O handler de imagens atua apenas sobre placeholders resolvidos, dissociando a comunicação de rede de alta latência (uploads de arquivos) do raciocínio cognitivo do modelo de linguagem.
\end{itemize}
Essa modularidade assegura escalabilidade e facilidade de manutenção no ciclo de vida de desenvolvimento do sistema.

\section{Metodologia de Validação e Testes}
Para garantir a estabilidade e a corretude funcional do agente, implementamos testes unitários e de integração sob o framework \lstinline{pytest} executados de forma assíncrona.

\subsection{Testes Unitários}
O arquivo \lstinline{tests/test_image_handler.py} valida de forma isolada a rotina de upload e substituição de tags. Utilizando dublês de teste (\lstinline{unittest.mock.patch}), interceptamos as chamadas de rede para a API de mídias e simulamos as respostas de upload de arquivos com URLs fictícias. O teste valida se os marcadores textuais foram integralmente convertidos para tags HTML válidas e se os metadados de arquivos foram anexados adequadamente nos locais corretos do JSON resultante.

\subsection{Testes de Integração}
Os arquivos \lstinline{tests/test_api_v3.py} e \lstinline{tests/test_ingestion_api.py} cobrem as rotas expostas pelo FastAPI. Simulamos um diretório com arquivos e testamos a chamada ao endpoint \lstinline{/ingest-question-bank} usando um cliente HTTP de testes assíncronos (\lstinline{httpx.AsyncClient}). Os testes confirmam que a tarefa em segundo plano (\lstinline{BackgroundTasks}) é enfileirada, que os arquivos de logs de auditoria são preenchidos corretamente e que o endpoint de status reflete a progressão dos arquivos e das questões de forma precisa.

\section{Perguntas de Pesquisa e Métricas de Avaliação}
Para quantificar a eficácia dos Large Language Models como assistentes e avaliar a robustez da arquitetura gerada para a plataforma PLURI, definimos as seguintes Questões de Pesquisa (RQs) em alinhamento com a proposta de mestrado do projeto:
\begin{itemize}
\item \textbf{RQ1 (Eficácia na Geração de Código):} Como diferentes modelos (Gemini 1.5 Pro vs. Flash) se comportam na geração de módulos de parsing e integração que atendam a requisitos de negócio reais?
\item \textbf{RQ2 (Design e Modelagem):} Em que medida o uso de LLMs facilita a definição de contratos de dados (schemas Pydantic) e a estruturação de uma arquitetura modular?
\item \textbf{RQ3 (Qualidade e Refatoração):} Qual o nível de intervenção humana e refatoração necessário para tornar o código gerado pelo LLM "pronto para produção"?
\item \textbf{RQ4 (Técnicas de Prompting):} Como a transição de Zero-shot para Few-shot impacta a robustez do código de extração gerado?
\end{itemize}

Para responder de forma empírica a essas perguntas, empregamos as seguintes métricas de Engenharia de Software:
\begin{enumerate}
\item \textbf{Aderência ao Requisito (Functional Correctness):} Mede o percentual de funções e rotinas de pipeline geradas de forma assistida que passam integralmente na suíte de testes unitários e de integração com a API da plataforma PLURI.
\item \textbf{Esforço de Refatoração (Edit Distance):} Medição da distância de edição (quantidade de adições, remoções e alterações manuais de linhas) necessária para estabilizar o código gerado pelo LLM e torná-lo aderente ao ambiente de produção.
\item \textbf{Complexidade Ciclomática:} Medida qualitativa e quantitativa da simplicidade, legibilidade e facilidade de manutenção estrutural do código de parsing gerado (comparando a saída do Gemini Pro versus o modelo Flash).
\item \textbf{Densidade de Erros de Integração:} Quantidade de quebras de contrato de payload (erros de validação Pydantic e falhas Beans HTTP 400/500 detectadas do lado do servidor) ocorridas ao longo das iterações de desenvolvimento do agente.
\end{enumerate}

\section{Análise de Custos, Desempenho e Engenharia de Integração}
Os custos de execução foram monitorados em tempo real. A precificação é baseada no volume de tokens de entrada e saída. A fórmula de custo total utilizada pelo pipeline é definida por:

\begin{equation}
\text{Custo Total} = \left( \frac{\text{Tokens In}}{10^6} \times P_{\text{in}} \right) + \left( \frac{\text{Tokens Out}}{10^6} \times P_{\text{out}} \right)
\end{equation}

Onde $P_{\text{in}}$ e $P_{\text{out}}$ representam os preços por milhão de tokens definidos no arquivo de configuração do sistema para cada modelo comercial disponível no Google AI Studio / Vertex AI:

\begin{table}[h]
\centering
\begin{tabular}{|l|c|c|}
\hline
\textbf{Modelo} & \textbf{Preço Input ($P_{\text{in}}$)} & \textbf{Preço Output ($P_{\text{out}}$)} \\ \hline
Gemini 1.5 Flash & \$0.075 & \$0.30 \\ \hline
Gemini 2.0 Flash & \$0.100 & \$0.40 \\ \hline
Gemini 2.5 Flash Lite & \$0.075 & \$0.30 \\ \hline
Gemini 1.5 Pro & \$1.250 & \$5.00 \\ \hline
\end{tabular}
\caption{Tabela de preços por 1M de tokens (USD).}
\label{tab:precos}
\end{table}

\subsection{Discussão de Trade-off Econômico}
Durante o processamento do lote do banco de questões (cerca de 100 arquivos contendo aproximadamente 400 questões no total), avaliamos o custo operacional aproximado de ingestão:
- **Abordagem Híbrida (Flash Lite para Detecção + Pro para Extração):** Custou cerca de \$2.45 no total. Obteve acurácia de 96\% no mapeamento de assuntos de eletrônica.
- **Abordagem Pura (Gemini 2.5 Flash Lite):** Custou aproximadamente \$0.18 no total. Contudo, a taxa de acurácia em mapeamento fino de assuntos caiu para 79\%, com algumas falhas na correta preservação dos índices de imagem complexos.
Desta forma, para produção, a arquitetura modular permite parametrizar dinamicamente o modelo ideal de acordo com a criticidade de cada documento técnico.

\subsection{Code Smells Identificados e Engenharia de Refatoração}
Durante as sessões de validação e testes de integração de ponta a ponta com a API do backend PLURI, nos deparamos com falhas recorrentes no código de validação e persistência (JPA) do servidor Java. A análise do código revelou a presença de diversos \textit{code smells} (maus cheiros no código) clássicos, conforme classificados por Fowler \cite{fowler1999refactoring}, os quais exigiram engenharia de refatoração sistemática:

\begin{enumerate}
\item \textbf{Incompatibilidade Tipo-Restrição (Constraint-Type Mismatch) [Smell: Primitive Obsession]:} O DTO \lstinline{DadosCriarQuestao.java} continha a anotação \lstinline{@NotEmpty} no campo \lstinline{Long area}. Em frameworks Jakarta Bean Validation, \lstinline{@NotEmpty} é incompatível com tipos numéricos como \lstinline{Long}, gerando a exceção \lstinline{HV000030} em tempo de execução. O \textit{smell} foi mitigado refatorando-se o campo para utilizar a validação correta de nulidade \lstinline{@NotNull}.
\item \textbf{Acoplamento Temporal e Efeitos Colaterais de Persistência (Cascade Persistence Bug) [Smell: Temporary Field / Side Effects]:} O método de cadastro de questões no servidor Java persistia individualmente as entidades de arquivos recebidas no payload JSON da requisição e, subsequentemente, tentava salvar a entidade \lstinline{Questao} correspondente (que possuía mapeamento JPA de cascata \lstinline{cascade = CascadeType.ALL}). Essa redundância temporal fazia com que o Hibernate tentasse registrar duas vezes os mesmos objetos com IDs pré-existentes, gerando a exceção fatal \lstinline{detached entity passed to persist}. O acoplamento foi refatorado no cliente Python, limpando-se o array \lstinline{arquivos} no payload JSON de envio. A persistência de mídia foi delegada de forma limpa ao método interno do servidor \lstinline{processarImagensDoHTML}, que varre o corpo HTML de forma isolada, gerando registros livres de conflito.
\item \textbf{Risco de Unboxing Implícito e Inicialização Nula [Smell: Incomplete Library Class / Missing Defaults]:} O campo \lstinline{rascunho} da classe de entidade \lstinline{Questao.java} era declarado como um wrapper \lstinline{Boolean} sem inicialização padrão (iniciando em \lstinline{null}). Na conversão de saída para o DTO \lstinline{DadosDetalhamentoQuestao.java}, que define a propriedade como tipo primitivo \lstinline{boolean}, o Java executava um auto-unboxing implícito (\lstinline{null.booleanValue()}), disparando uma \lstinline{NullPointerException}. Corrigimos o código aplicando a inicialização explícita \lstinline{private Boolean rascunho = false;} diretamente na entidade JPA.
\item \textbf{Ausência de Tratamento de Exceções de Domínio (Semantic Missing Data):} A ausência de assuntos identificados pelo LLM resultava em erros de validação direta de campos vazios na API. Implementamos uma rotina de fallback inteligente no cliente Python, injetando IDs de assuntos válidos derivados da área pedagógica sob análise antes do despacho do JSON.
\end{enumerate}

\section{Conclusão}
A modernização do Agente PLURI com rotinas inteligentes de desduplicação e mapeamento de imagens, ingestão em lote, tratamento de quotas limite da Vertex AI via Sliding Window e a resolução sistemática de conflitos de arquitetura e validação de dados com o ecossistema Spring Boot/JPA solidifica o valor de assistentes baseados em LLMs em pipelines de software corporativos. A mitigação de falhas arquiteturais críticas foi alcançada através de uma especificação contratual rígida combinada a um processamento estruturado e limpo.

Trabalhos futuros focarão na automação de processos de correção autônoma de falhas de envio da API (self-healing), onde o próprio agente interpreta respostas HTTP 400/500 da plataforma PLURI e refatora dinamicamente o JSON enviado utilizando o LLM antes de efetuar novas tentativas de cadastro.

\bibliographystyle{ACM-Reference-Format}
\bibliography{references.bib}

\end{document}
\endinput
